\subsection{Occupancy and boundary-to-bulk transgression}

With the family bridge installed, the matter space becomes $S^+\otimes F$ and the occupancy $\Oadm=48$ is already fixed by the preceding proposition. Here \emph{occupancy} means the admissible chiral-state count after family lifting; it is not used as a primitive datum. The bridge point is not that family lifting creates a new numerical surcharge. Rather, the inherited hard kernel value
\[
\deltatop=\frac{3}{256\pi^4}
\]
admits on the canonical family branch the occupancy reinterpretation $\deltatop=\Oadm\,\cthree^4$.
Likewise the hard second-order defect contribution is re-read as
\[
\delta_2=\frac54\,\deltatop^2=B\gamma\,\Oadm^2\cthree^8.
\]
This bridge reinterpretation is the boundary-to-bulk step that later feeds both the electromagnetic and geometric benchmark channels. The coefficient $5/4$ is therefore not an extra numerical motif; it is the minimal-carrier compression factor $B\gamma$ rewritten in defect language. At the same time, the equality $\deltatop=\Oadm\cthree^4$ should be read exactly as a bridge reinterpretation, not yet as a proved transgression theorem in the ordinary field-theory sense.

\subsection{From bridge to geometric completion}

\begin{principle}[Completion split]
Once the carrier packet and family bridge are fixed, the focused main-text completion problem splits into a geometric program $\Ageom$ and a transport program $\Atrans$. Record algebra is retained only as a conditional appendix layer in \Cref{app:record}.
\end{principle}

The two main-text programs are packaged as
\[
\Ageom
:=
\langle E_3\oplus E_2,\ \chi,\ \mathcal B_{\mathrm{rel}},\ g,\ \omega_{\mathrm{spin}}\rangle,
\qquad
\Atrans
:=
\langle U_6,\ T,\ D_y,\ \mathcal Y_y^{(\varepsilon)},\ \mathcal C_{\Seam}\rangle.
\]
This split keeps the carrier theorem, the geometric closure problem, and the record-level extensions from being flattened into one undifferentiated claim.

\subsection{Relative APS and superconnection setup}

The geometric completion is formulated on a graded bundle
\[
\mathcal E_{\mathrm{rel}}=\mathcal E_{\mathrm{rel}}^+\oplus\mathcal E_{\mathrm{rel}}^-,
\]
with a declared reference pair $(D,D_{\mathrm{ref}})$ and APS-type boundary conditions on the seam $\Seam$; see Appendix~\ref{app:aps}. The geometric package is
\[
\mathbb A_{\Seam}:=D_{\mathrm{rel}}\oplus\mathcal B_{\mathrm{rel}}\oplus\nabla_\chi.
\]
This is a \emph{construction}, not yet a theorem. Its value is to keep grading, reference subtraction, seam data, and the $\chi$ channel in one object from the start rather than mixing them only at the level of later heuristics.

\subsection{Bosonic relative index and Higgs selection}

\begin{conjecture}[Bosonic relative index closure]
The bosonic relative index satisfies
\[
\Ind_{\mathrm{rel}}^{\mathsf{Adm}}(\mathcal B_{E_2}^+)=1,
\qquad
\Ind_{\mathrm{rel}}^{\mathsf{Adm}}(\mathcal B_{E_3}^+)=0.
\]
Consequently the unique light seam-even bosonic zero mode lies in
\[
\Phi\in\Gamma(E_2)\cong (1,2)_{1/2},
\qquad
\inv\Phi=+\Phi.
\]
\end{conjecture}

This is the sharp form of the Higgs claim used here. The issue is not whether a Higgs doublet can be written down, but whether the relative bosonic index can be shown to select it uniquely.
If that selection theorem closes, the renormalizable carrier-compatible Yukawa couplings are the standard ones,
\[
Q\,u^c\,\Phi,
\qquad
Q\,d^c\,\Phi^\dagger,
\qquad
L\,e^c\,\Phi^\dagger,
\qquad
L\,\nu^c\,\Phi,
\]
so the Higgs statement becomes a genuine structural bridge rather than a notation choice.

\subsection{Heat-kernel coefficient map and \texorpdfstring{$\chi$}{chi} bifurcation}

The local target action is written as
\begin{equation}
\label{eq:chi-effective-action}
\begin{aligned}
\Gammaeff
=
\int d^4x\sqrt{-g}\,
\Big[
c_R\chi^2R
+Z_\chi(\partial\chi)^2
-\mu_\Phi^2(\chi)\,\Phi^\dagger\Phi
+\lambda_\Phi(\Phi^\dagger\Phi)^2
+\frac{1}{6M^2}R^2
+\cdots
\Big].
\end{aligned}
\end{equation}
The coefficient bookkeeping expected from the relative spectral expansion is summarized in \Cref{tab:heat-kernel-map}.

\begin{table}[t]
\centering
\small
\begin{tabularx}{\linewidth}{@{}>{\raggedright\arraybackslash}p{0.15\linewidth}>{\raggedright\arraybackslash}p{0.22\linewidth}>{\raggedright\arraybackslash}p{0.22\linewidth}>{\raggedright\arraybackslash}p{0.14\linewidth}>{\raggedright\arraybackslash}p{0.14\linewidth}@{}}
\toprule
Coefficient & Operator sector & Generated term & $\chi$ scaling & Status \\
\midrule
$a_0$ & relative bulk trace & vacuum-energy bookkeeping & $\chi^4$ & Hard \\
$a_2$ & even bulk term & $\chi^2R$, mass-type terms & $\chi^2$ & Completion \\
$a_4$ & local curvature and scalar sector & $(\partial\chi)^2$, $\Phi^\dagger\Phi$, $(\Phi^\dagger\Phi)^2$, $R^2$ & $\chi^0$ & Completion \\
boundary/defect term & seam correction and parity data & local defect counterterms, seam-even selection & $\Delta_{\Seam}$ & Bridge \\
$\eta$ term & APS spectral asymmetry & sign and spectral-flow control & scale free & Bridge \\
\bottomrule
\end{tabularx}
\caption{Heat-kernel map for the geometric completion program.}
\label{tab:heat-kernel-map}
\end{table}

The resulting geometric logic is displayed in \Cref{fig:chi-bifurcation}. The same scale channel is asked to speak both to the induced Einstein term and to the scalar vacuum response. This is why the $\chi$ sector is the principal unresolved geometric closure problem.

\begin{figure}[t]
\centering
\begin{tikzpicture}[
    >=Latex,
    node distance=1.8cm and 2.1cm,
    box/.style={draw=orange!70!black, rounded corners, fill=orange!5, align=center, minimum width=2.4cm, minimum height=0.95cm}
]
\node[box] (chi) {$\chi$ channel};
\node[box, below left=of chi] (einstein) {$\Mpl^2\sim \chi_0^2$\\Einstein branch};
\node[box, below=of chi] (higgs) {$v^2=\mu_\Phi^2(\chi_0)/\lambda_\Phi$\\Higgs branch};
\node[box, below right=of chi] (r2) {$R^2/M^2$\\higher-curvature branch};
\draw[->, thick] (chi) -- (einstein);
\draw[->, thick] (chi) -- (higgs);
\draw[->, thick] (chi) -- (r2);
\end{tikzpicture}
\caption{Bifurcation of the $\chi$ channel. The three arrows are treated as one coupled geometric completion problem rather than as independent numerical inputs.}
\label{fig:chi-bifurcation}
\end{figure}

\begin{conjecture}[Induced Einstein and vacuum channel]
On the canonical branch the relative heat-kernel closure yields
\begin{equation}
\label{eq:induced-einstein}
\Mpl^2=\frac{\chi_0^2}{2\pi^2},
\qquad
v^2=\frac{\mu_\Phi^2(\chi_0)}{\lambda_\Phi}.
\end{equation}
\end{conjecture}

\begin{remark}[Newton form is not an independent output]
Using the standard reduced-Planck definition $\Mpl^{-2}=8\pi G_N$, the first relation in \Cref{eq:induced-einstein} is equivalently
\[
G_N\chi_0^2=\frac{\pi}{4}.
\]
This Newton-form equation is therefore a rewriting of the Einstein-branch closure, not a second independent prediction.
\end{remark}

\begin{conjecture}[Carrier-dressed electroweak vacuum]
If the $\chi$-channel Higgs mass parameter inherits the carrier rank and transport seed, the sharpest conditional target formula is
\begin{equation}
\label{eq:carrier-dressed-v}
v
=
\Mpl\,\gcar\,\beta_{\mathrm{rad}}^2\,
\exp\!\left(-\frac{\alpha^{-1}(0)}{5}\right)
\exp\!\left(-\frac{\delta_{\mathrm{ph}}}{5}\right),
\end{equation}
where $\gcar=5$ is the carrier rank and $\delta_{\mathrm{ph}}=\frac35+\frac{\phiz}{6}$ is the transport bridge offset defined in \Cref{eq:transport-bridge-offset}. This makes the electroweak vacuum a carrier-rank-weighted, $\alpha$-suppressed, transport-corrected fraction of the Planck scale, and turns the v2.8 constant-factory observation into a theorem candidate once the $a_2$/$a_4$ heat-kernel coefficients are computed in a declared scheme.
\end{conjecture}

\begin{remark}[Sector-dressed \texorpdfstring{$\chi$}{chi} channels]
The quantity $\chi_0$ entering the induced Einstein channel is not identified with the transport-sector confinement scale. In the hadronic sector we write $\chi_{\QCD}$ and regard it as a locally dressed channel obtained from the same underlying $\chi$ response after sector-dependent renormalization and confinement-scale matching:
\[
\chi \longrightarrow \chi_0 \quad \text{(gravitational background)},
\qquad
\chi \longrightarrow \chi_{\QCD} \quad \text{(transport-dressed local channel)}.
\]
This is the minimal statement needed to avoid the otherwise immediate scale clash between the Planck bridge and the QCD string tension.
\end{remark}

The E$_8$ scale ladder is not used as a primitive input in the main text. Only its short downstream grammar is retained in the dedicated main-text section, while the fuller arithmetic continuations are deferred to Appendix~\ref{app:constants}.

\section{Completion program B: Yukawas, masses, flavor, confinement, and strong CP}

The second completion block is the transport compiler. Its role is to turn the carrier packet and the family/occupancy bridge into actual masses, mixings, hadronic admissibility, and strong-CP control. The logic is again staged: first the discrete phase operator and the transport kernel, then path-length masses and mixing data, then confinement and strong CP as transport-side admissibility statements.

\subsection{Discrete phase operator and transport grammar}

Let
\[
\omega=e^{2\pi i/3},
\qquad
U_6=\mathrm{diag}(1,\omega,\omega^2,-1,-\omega,-\omega^2).
\]
This makes the phase grammar concrete rather than rhetorical. The six admissible phases are shown in \Cref{fig:z6-circle}, and CP acts by inversion $\delta\mapsto-\delta$ on this circle.

\begin{remark}[Hexagonal spectral geometry]
The operator $U_6$ is a unitary circulant with eigenphases $e^{in\pi/3}$ for $n=0,\dots,5$. The transport kernel $D_y=y\,\mathbf 1-\delta_{\mathrm{ph}}U_6$ is therefore a spectral problem on the discrete $C_6$ cycle. CKM and PMNS mixings, mass hierarchies, and the Möbius quotients of \Cref{eq:mobius-quotient} are then not loose analytic fits but spectral-theory statements on a regular hexagon. The near-pole hierarchy at $y\approx\delta_{\mathrm{ph}}$ becomes the resonance of a hexagonal resolvent rather than tuned parameter choice.
\end{remark}

\begin{figure}[t]
\centering
\begin{tikzpicture}[scale=1.2,>=Latex]
\draw[thick] (0,0) circle (2.0);
\foreach \ang/\lab in {90/{0},30/{\pi/3},-30/{2\pi/3},-90/{\pi},-150/{4\pi/3},150/{5\pi/3}} {
    \filldraw[blue!60!black] ({2*cos(\ang)},{2*sin(\ang)}) circle (2pt);
    \node at ({2.45*cos(\ang)},{2.45*sin(\ang)}) {$\lab$};
}
\draw[->,red!70!black,thick] (1.25,0.9) arc[start angle=35,end angle=-35,radius=1.55];
\node[red!70!black] at (2.45,0.1) {\small CP};
\end{tikzpicture}
\caption{$\mathbb Z_6$ phase circle with CP inversion.}
\label{fig:z6-circle}
\end{figure}

\subsection{Regularized Yukawa kernels and path-length masses}

Here
\begin{equation}
\label{eq:transport-bridge-offset}
\delta_{\mathrm{ph}}:=\frac35+\frac{\phiz}{6}
\end{equation}
is treated as a transport bridge parameter inherited from the cross-ratio and boundary-monodromy program, rather than as a proved hard-kernel identity. The full derivation is therefore deferred to the supplementary layer, while the regularized transport operator used here is
\begin{equation}
\label{eq:yukawa-kernel}
\begin{aligned}
D_y&=y\,\mathbf 1-\delta_{\mathrm{ph}}U_6,\\
\mathcal Y_y^{(\varepsilon)}&=(D_y^\dagger D_y+\varepsilon^2)^{-1},
\end{aligned}
\end{equation}
and the singular limit $\varepsilon\to0$ is taken only in admissible sectors.

Because $U_6$ has eigenphases $e^{in\pi/3}$, the regularized transport spectrum is
\begin{equation}
\label{eq:transport-spectrum}
\lambda_{y,n}^{-1}
=
y^2+\delta_{\mathrm{ph}}^2-2y\delta_{\mathrm{ph}}\cos\!\left(\frac{n\pi}{3}\right)+\varepsilon^2,
\qquad
n=0,\dots,5,
\end{equation}
for the eigenvalues $\lambda_{y,n}$ of $\mathcal Y_y^{(\varepsilon)}$. In particular,
\begin{equation}
\label{eq:mobius-quotient}
\frac{\lambda_{y,0}}{\lambda_{y,3}}
=
\left(\frac{y+\delta_{\mathrm{ph}}}{y-\delta_{\mathrm{ph}}}\right)^2.
\end{equation}
\begin{remark}[Near-pole hierarchy mechanism]
On the retained branch
\[
\delta_{\mathrm{ph}}=\frac35+\frac{\phiz}{6}\approx \num{0.60886}.
\]
Hence the Möbius quotient contains a built-in near-pole whenever $y$ approaches $\delta_{\mathrm{ph}}$, making large hierarchy ratios structurally plausible. At present this is a hierarchy mechanism, not yet a theorem of explanatory control: until admissible path lengths $L_{f,j}$ and bounded residual factors $\Lambda_{f,j}$ are derived geometrically, the same near-pole can also masquerade as a tuning surface.
\end{remark}

Absolute masses are attached to admissible path lengths,
\begin{equation}
\label{eq:path-length-yukawa}
y_{f,j}^{\mathrm{UV}}=\lambdaY^{L_{f,j}}\Lambda_{f,j},
\qquad
\lambdaY=\sqrt{\phiz(1-\phiz)}.
\end{equation}
After the electroweak bridge fixes $v$, the physical mass rows are emitted through
\begin{equation}
\label{eq:mass-bridge}
m_{f,j}=\frac{y_{f,j}v}{\sqrt2}.
\end{equation}
The role of the residual factors $\Lambda_{f,j}$ must be kept explicit. They are bounded $O(1)$ transport prefactors, not unrestricted fit knobs.

\begin{remark}[Flavor cusps as carrier cusps]
The admissible cusp values $y\in\{1,\frac13,\frac23\}$ are not arbitrary rational points; they are the absolute values of the singular hypercharges $|Y(e^c)|=1$, $|Y(d^c)|=\frac13$, $|Y(u^c)|=\frac23$. The transport grammar therefore lives on the same rational grid as the carrier packet itself.
\end{remark}

On the minimal admissible branch the word-length candidates are
\begin{equation}
\label{eq:minimal-word-lengths}
\begin{gathered}
L_t=0,\quad
L_b=L_\tau=3,\quad
L_c=4,\quad
L_s=L_\mu=5,\\
L_d=7,\quad
L_u=L_e=8.
\end{gathered}
\end{equation}
These are not free parameters; they are the minimal admissible word lengths on the $C_6$ graph that reproduce the observed mass hierarchy through $\lambdaY^{L_{f,j}}$. The assignment is unique up to residual-factor control: once the hierarchy ordering is fixed, the nearest-integer graph distance is determined by $\lambdaY\approx\num{0.2244}$.

The transport-side status is summarized in \Cref{tab:transport-dictionary}.

\begin{table}[t]
\centering
\scriptsize
\resizebox{\linewidth}{!}{%
\begin{tabular}{@{}lllll@{}}
\toprule
Sector & Admissible length data & Residual prefactor & Role & Status \\
\midrule
charged leptons & $L_{e}=8,\;L_\mu=5,\;L_\tau=3$ & $\Lambda_e,\Lambda_\mu,\Lambda_\tau$ & ladder and hierarchy & bounded; operator closure open \\
up quarks & $L_u=8,\;L_c=4,\;L_t=0$ & $\Lambda_u,\Lambda_c,\Lambda_t$ & up-type ratio ladder & bounded; phase closure open \\
down quarks & $L_d=7,\;L_s=5,\;L_b=3$ & $\Lambda_d,\Lambda_s,\Lambda_b$ & down-type ratio ladder & bounded; phase closure open \\
quark mixing & $L_{ij}^{q}$ & $\Lambda_{ij}^{q}$ & CKM transport entries & bounded; phase closure open \\
neutrino mixing & $L_{ij}^{\nu}$ & $\Lambda_{ij}^{\nu}$ & PMNS transport entries & bounded; bridge partial \\
hadron basis & center-neutral word lengths & block coefficients & meson and baryon transport basis & open bound-state layer \\
strong CP & determinant spectrum of $\mathcal Y_f$ & none if positivity closes & phase suppression & positivity hardening pending \\
\bottomrule
\end{tabular}}
\caption{Transport algebra dictionary with minimal admissible word-length candidates.}
\label{tab:transport-dictionary}
\end{table}

\subsection{CKM and PMNS sharpening}

The leading CKM magnitudes on the transport branch are
\begin{equation}
\label{eq:ckm-magnitudes}
s_{12}=\lambda_C,
\qquad
s_{23}=\frac{\phiz}{1+\lambda_C},
\qquad
s_{13}=\frac{\lambda_C^3}{3}.
\end{equation}
The CKM phase receives a tri-familial dressing beyond the raw topological $\pi/3$:
\begin{equation}
\label{eq:ckm-phase}
\delta_{\mathrm{CKM}}=\frac{\pi}{3}+3\lambda_C^2\approx\num{1.198}\;\mathrm{rad},
\end{equation}
so $\pi/3$ remains the base monodromy phase and $3\lambda_C^2$ is the leading family-transport correction.

For the neutrino sector the minimal seesaw closure on the transport branch is
\begin{equation}
\label{eq:neutrino-masses}
m_3=e^{2\lambda_C^2}\frac{v^2}{M_R},
\qquad
m_2=\pi\phiz\,m_3,
\qquad
m_1=0,
\end{equation}
with $M_R\approx\num{1.31e15}\;\mathrm{GeV}$ fixed by the seesaw E$_8$ block at $(n,r,I_1)=(5,5,\frac{1}{12})$. Numerically this gives
\[
m_3\approx\num{0.0508}\;\mathrm{eV},
\qquad
m_2\approx\num{0.0085}\;\mathrm{eV},
\]
consistent with normal ordering.

\subsection{Confinement as admissible word selection}

\begin{proposition}[Admissible color words]
For each carrier word $w$, let $z(w)\in\{1,\omega,\omega^2\}$ be its color-center charge and let $\varepsilon(w)\in\{\pm1\}$ be its seam parity. Physical hadronic words satisfy
\[
z(w)=1,
\qquad
\varepsilon(w)=+1.
\]
\end{proposition}

The confinement claim is thus stated as an admissibility rule, not as a late dynamical slogan. A schematic version is shown in \Cref{fig:admissible-words}.

\begin{figure}[t]
\centering
\begin{tikzpicture}[>=Latex, node distance=1.1cm and 1.3cm]
\node[draw=red!70!black, fill=red!4, rounded corners, minimum width=2.1cm, minimum height=1.0cm, align=center] (q) {free quark\\$z(w)\neq 1$};
\node[draw=green!50!black, fill=green!4, rounded corners, minimum width=2.1cm, minimum height=1.0cm, align=center, right=of q] (m) {meson\\$z(w)=1$};
\node[draw=green!50!black, fill=green!4, rounded corners, minimum width=2.1cm, minimum height=1.0cm, align=center, right=of m] (b) {baryon\\$z(w)=1$};
\draw[red!70!black, very thick] ($(q.north west)+(0.12,-0.12)$) -- ($(q.south east)+(-0.12,0.12)$);
\draw[red!70!black, very thick] ($(q.south west)+(0.12,0.12)$) -- ($(q.north east)+(-0.12,-0.12)$);
\node[below=0.2cm of q] {\small excluded};
\node[below=0.2cm of m] {\small admissible};
\node[below=0.2cm of b] {\small admissible};
\end{tikzpicture}
\caption{Confinement as admissible word selection.}
\label{fig:admissible-words}
\end{figure}

The transport-sector string tension is written as
\[
\sigma=\cthree^2\chi_{\QCD}^2,
\]
so the relative Wilson loop obeys
\[
\langle W(C)\rangle_{\mathrm{rel}}\sim e^{-\sigma A(C)}.
\]
The sector-dependent scale $\chi_{\QCD}$ is the minimal notation needed to avoid confusing the hadronic transport channel with the gravitational $\chi_0$ background; the two channels are linked by the common underlying $\chi$ response but not by a naive equality of scales.

